% !TeX root = ../cosmology_summary.tex

%% --------------------------------------------------------
\chapter{Практичний посібник: від даних CMB до параметрів моделі}
\label{sec:practical}
%% --------------------------------------------------------

Цей додаток відповідає на одне питання: як людина, що прочитала
цей конспект, може взяти реальні дані Planck і власноруч
витягти з них космологічні параметри? Нижче~--- покроковий
маршрут від встановлення пакетів до трикутного графіка
posterior-розподілів.

Маршрут складається з чотирьох кроків:
\begin{enumerate}
  \item Порахувати теоретичний спектр $C_\ell^{TT}$ за допомогою CAMB.
  \item Завантажити дані Planck і реалізувати простий $\chi^2$-фіт вручну.
  \item Підключити байєсівський семплер Cobaya і запустити MCMC.
  \item Прочитати результат: posterior-розподіли і довірчі контури.
\end{enumerate}

%% --------------------------------------------------------
\section{Крок 1. Теоретичний спектр через CAMB}
%% --------------------------------------------------------

Встановлення:
\begin{minted}{bash}
pip install camb
\end{minted}

Мінімальний скрипт, який рахує $\mathcal{D}_\ell^{TT}
= \ell(\ell+1)C_\ell^{TT}/(2\pi)$ для стандартної моделі
$\Lambda$CDM з параметрами Planck~2018:

\begin{minted}{python}
import camb
import numpy as np
import matplotlib.pyplot as plt

# --- Задаємо параметри моделі ---
pars = camb.CAMBparams()
pars.set_cosmology(
    H0=67.36,          # параметр Хаббла [км/с/Мпк]
    ombh2=0.02237,     # фізична густина баріонів Omega_b h^2
    omch2=0.1200,      # фізична густина темної матерії Omega_c h^2
    tau=0.0544,        # оптична глибина реіонізації
)
pars.InitPower.set_params(
    As=2.1e-9,         # амплітуда первинних збурень
    ns=0.9649,         # спектральний індекс
)
pars.set_for_lmax(2500)

# --- Рахуємо спектр ---
results = camb.get_results(pars)
powers  = results.get_cmb_power_spectra(pars, CMB_unit='muK')
Dl_TT   = powers['total'][:, 0]   # D_ell^TT [мкК^2]
ells    = np.arange(len(Dl_TT))

# --- Малюємо ---
plt.figure(figsize=(8, 4))
plt.plot(ells[2:], Dl_TT[2:])
plt.xlabel(r'$\ell$')
plt.ylabel(r'$\mathcal{D}_\ell^{TT}\ [\mu\mathrm{K}^2]$')
plt.xscale('log')
plt.tight_layout()
plt.savefig('spectrum_lcdm.pdf')
\end{minted}

\noindent
Якщо все встановлено правильно, ви отримаєте графік з трьома
добре видимими акустичними піками~--- той самий спектр, що
і на рис.~\ref{fig:cmb_tt}.

\textit{Зауваження про одиниці.} CAMB повертає
$\mathcal{D}_\ell = \ell(\ell+1)C_\ell/(2\pi)$ у мкК$^2$.
Саме в цих одиницях публікуються дані Planck, тому порівняння
пряме.

%% --------------------------------------------------------
\section{Крок 2. Дані Planck і простий $\chi^2$-фіт}
%% --------------------------------------------------------

\subsection{Завантаження даних}

Офіційні дані Planck~2018 доступні на сервері \href{https://pla.esac.esa.int/}{ESA}.
Для початку достатньо бінованих TT-даних з файлу
\texttt{COM\_PowerSpect\_CMB-TT-binned\_R3.01.txt}
(колонки: $\ell$, $\mathcal{D}_\ell$, $\sigma^-$, $\sigma^+$).

\begin{minted}{python}
import numpy as np

data = np.loadtxt('COM_PowerSpect_CMB-TT-binned_R3.01.txt')
ell_data  = data[:, 0]
Dl_data   = data[:, 1]
sigma_m   = data[:, 2]   # нижня похибка
sigma_p   = data[:, 3]   # верхня похибка
sigma     = 0.5 * (sigma_m + sigma_p)   # симетризована похибка
\end{minted}

\subsection{Проста $\chi^2$-статистика}

Щоб зрозуміти принцип, реалізуємо наївний $\chi^2$-фіт вручну~---
без likelihood Planck, але з правильною фізичною логікою:
\begin{equation}\label{eq:chi2}
  \chi^2(\theta) = \sum_\ell
    \frac{\bigl[\mathcal{D}_\ell^{\rm data}
               - \mathcal{D}_\ell^{\rm theory}(\theta)\bigr]^2}
         {\sigma_\ell^2},
\end{equation}
де $\theta = \{H_0, \Omega_b h^2, \Omega_c h^2, \tau, A_s, n_s\}$.

\begin{minted}{python}
from scipy.optimize import minimize
from scipy.interpolate import interp1d

def get_theory(H0, ombh2, omch2, tau, As, ns):
    """Повертає D_ell^TT на точках ell_data."""
    pars = camb.CAMBparams()
    pars.set_cosmology(H0=H0, ombh2=ombh2,
                       omch2=omch2, tau=tau)
    pars.InitPower.set_params(As=As, ns=ns)
    pars.set_for_lmax(int(ell_data.max()) + 50)
    results = camb.get_results(pars)
    powers  = results.get_cmb_power_spectra(
                  pars, CMB_unit='muK')
    Dl_full = powers['total'][:, 0]
    ells_full = np.arange(len(Dl_full))
    interp = interp1d(ells_full, Dl_full,
                      kind='linear', fill_value='extrapolate')
    return interp(ell_data)

def chi2(params):
    H0, ombh2, omch2, tau, log_As, ns = params
    As = np.exp(log_As) * 1e-10
    Dl_th = get_theory(H0, ombh2, omch2, tau, As, ns)
    return np.sum(((Dl_data - Dl_th) / sigma)**2)

# Стартова точка --- значення Planck 2018
x0 = [67.36, 0.02237, 0.1200, 0.0544, np.log(21.0), 0.9649]
result = minimize(chi2, x0, method='Nelder-Mead',
                  options={'maxiter': 500, 'xatol': 1e-3})
print("Best-fit params:", result.x)
print("chi2_min / N_dof:", result.fun / len(ell_data))
\end{minted}

\noindent
\textit{Застереження.} Цей $\chi^2$ є наближенням: справжня
likelihood Planck враховує некомутативність похибок і
міжмультипольні кореляції. Для публікаційних результатів
використовують офіційний пакет \texttt{clik} або
\texttt{candl}. Але для розуміння принципу цього достатньо.

%% --------------------------------------------------------
\section{Крок 3. MCMC через Cobaya}
%% --------------------------------------------------------

$\chi^2$-мінімізація дає одну точку в просторі параметрів.
MCMC дає повний posterior-розподіл~--- тобто відповідає
на питання «наскільки добре параметр визначений?»

\subsection{Що таке MCMC в двох словах}

Алгоритм Метрополіса--Гастінгса будує ланцюжок точок
$\theta_1, \theta_2, \ldots, \theta_N$ у просторі параметрів
так, що їхня щільність пропорційна posterior-розподілу
$P(\theta | {\rm data}) \propto \mathcal{L}({\rm data}|\theta)
\cdot \pi(\theta)$, де $\pi(\theta)$~--- prior (попереднє
знання про параметри). Після збіжності ланцюжка гістограма
$\theta_i$ і є posterior.

\subsection{Встановлення і конфігурація}

\begin{minted}{python}
pip install cobaya
cobaya-install camb --packages-path ./packages
\end{minted}

Конфігурація у форматі YAML (\texttt{lcdm\_fit.yaml}):

\begin{minted}{python}
likelihood:
  planck_2018_highl_plik.TT:
    path: ./packages/data/planck_2018

theory:
  camb:
    path: ./packages/code/CAMB

params:
  H0:
    prior: {min: 60, max: 80}
    ref:   {mean: 67.36, std: 0.5}
    latex: H_0
  ombh2:
    prior: {min: 0.018, max: 0.026}
    ref:   {mean: 0.02237, std: 0.0002}
    latex: \Omega_b h^2
  omch2:
    prior: {min: 0.10, max: 0.14}
    ref:   {mean: 0.1200, std: 0.001}
    latex: \Omega_c h^2
  tau:
    prior: {min: 0.01, max: 0.12}
    ref:   {mean: 0.0544, std: 0.007}
    latex: \tau
  logA:
    prior: {min: 2.5, max: 3.5}
    ref:   {mean: 3.044, std: 0.014}
    drop:  true
    latex: \ln(10^{10}A_s)
  ns:
    prior: {min: 0.9, max: 1.05}
    ref:   {mean: 0.9649, std: 0.004}
    latex: n_s
  As:
    value: "lambda logA: 1e-10*np.exp(logA)"
    latex: A_s

sampler:
  mcmc:
    Rminus1_stop: 0.01    # критерій збіжності Gelman-Rubin
    max_samples:  10000

output: chains/lcdm
\end{minted}

Запуск:
\begin{minted}{python}
cobaya-run lcdm_fit.yaml
\end{minted}

\noindent
\textit{Критерій збіжності.} Параметр \texttt{Rminus1\_stop: 0.01}
означає критерій Гелмана--Рубіна $\hat{R} - 1 < 0.01$: якщо
кілька паралельних ланцюжків дають однаковий розподіл,
семплювання вважається збіжним. Типовий час на ноутбуці~---
кілька годин для TT-likelihood.

%% --------------------------------------------------------
\section{Крок 4. Читання результату}
%% --------------------------------------------------------

\subsection{Трикутний графік (corner plot)}

\begin{minted}{python}
pip install getdist

from getdist import plots, MCSamples
import getdist.plots as gdplt

# Зчитуємо chain
samples = MCSamples(
    root='chains/lcdm',
    settings={'ignore_rows': 0.3}   # відкидаємо burn-in
)

# Трикутний графік для двох параметрів
g = gdplt.get_subplot_plotter()
g.triangle_plot(
    [samples],
    ['H0', 'ns', 'ombh2'],
    filled=True,
    contour_colors=['steelblue']
)
g.export('posterior_triangle.pdf')
\end{minted}

Результат~--- матриця графіків: на діагоналі маргіналізовані
1D posterior-розподіли кожного параметра, поза діагоналлю~---
2D контури довіри (68\% і 95\%).

\subsection{Як читати posterior}

\begin{itemize}
  \item \textbf{Ширина 1D posterior}~--- це невизначеність
        параметра: чим вужча, тим краще дані його обмежують.
  \item \textbf{Нахил 2D контура}~--- кореляція між параметрами.
        Класичний приклад: $H_0$ і $\Omega_m$ антикорельовані
        через CMB~--- їхній добуток $\Omega_m h^2$ визначений
        краще ніж кожен окремо.
  \item \textbf{Зміщення центра}~--- якщо ваш best-fit далеко
        від таблиці Planck~2018, перевірте likelihood і prior.
\end{itemize}

\subsection{Похідні параметри: $\Omega_m$, $\sigma_8$, $S_8$}

MCMC семплює базові шість параметрів
$\{H_0, \Omega_b h^2, \Omega_c h^2, \tau, A_s, n_s\}$.
Але фізично цікавіші часто є \emph{похідні} параметри~---
функції базових. Cobaya вміє рахувати їх автоматично і
включати в chain.

\textbf{Параметри матерії.} Повна густину матерії
$\Omega_m = \Omega_b + \Omega_c$ виводиться з базових через:
\begin{equation}
  \Omega_m = \frac{\Omega_b h^2 + \Omega_c h^2}{(H_0/100)^2}.
\end{equation}
У конфігурації Cobaya:
\begin{minted}{yaml}
params:
  omegam:
    derived: >
      lambda ombh2, omch2, H0:
        (ombh2 + omch2) / (H0/100)**2
    latex: \Omega_m
  omegab:
    derived: "lambda ombh2, H0: ombh2 / (H0/100)**2"
    latex: \Omega_b
\end{minted}

\textbf{Параметр $\sigma_8$.} Стандартне відхилення
флуктуацій густини матерії на масштабі $8\,h^{-1}$~Мпк~---
це ключовий параметр для порівняння з даними слабкого
лінзування і підрахунку скупчень. CAMB рахує його сам:
\begin{minted}{yaml}
params:
  sigma8:
    derived: true   # CAMB повертає автоматично
    latex: \sigma_8
\end{minted}

\textbf{Параметр $S_8$.} Комбінація, яка найкраще
вимірюється слабким лінзуванням~--- вона «ортогональна»
до $H_0$ і тому є незалежним тестом моделі:
\begin{equation}
  S_8 \equiv \sigma_8\sqrt{\frac{\Omega_m}{0.3}}.
\end{equation}
\begin{minted}{yaml}
params:
  S8:
    derived: "lambda sigma8, omegam: sigma8*(omegam/0.3)**0.5"
    latex: S_8
\end{minted}

Після запуску MCMC всі ці параметри з'являться в chain
і getdist покаже їхні posterior-розподіли поряд з базовими.
Класичне питання: чи узгоджується $S_8$ з Planck і з даними
слабкого лінзування DES/KiDS? Саме ця \emph{$S_8$-tension}
($\sim 2$--$3\sigma$) є другою відкритою суперечністю
$\Lambda$CDM після Hubble tension.

\subsection{Модифікований спектр: зламаний первинний спектр}

Щоб перевірити модель зі зламом, достатньо замінити стандартний
первинний спектр у CAMB. Додаємо два нових параметри:
$k_{\rm break}$~--- положення зламу, $\Delta n$~--- стрибок
спектрального індексу:

\begin{minted}{python}
def broken_pk(k, As, ns, k_break, delta_n):
    """Первинний спектр зі зламом."""
    pk = As * (k / 0.05) ** (ns - 1)
    pk[k > k_break] *= (k[k > k_break] / k_break) ** delta_n
    return pk
\end{minted}

У конфігурації Cobaya додаємо нові параметри:
\begin{minted}{yaml}
params:
  k_break:
    prior: {min: 0.01, max: 1.0}   # [1/Мпк]
    ref:   {mean: 0.1, std: 0.02}
    latex: k_{\rm break}
  delta_n:
    prior: {min: -1.0, max: 1.0}
    ref:   {mean: 0.0, std: 0.1}
    latex: \Delta n
\end{minted}

Якщо posterior для $\Delta n$ зміститься від нуля більш ніж
на $2\sigma$~--- це свідчення на користь нестандартного спектра.

%% --------------------------------------------------------
\section{Типові помилки початківця}
%% --------------------------------------------------------

\begin{enumerate}
  \item \textbf{Забути про burn-in.} Перші $\sim 30\%$ кроків
        ланцюжка~--- це «прогрів», їх треба відкидати
        (\texttt{ignore\_rows: 0.3} у getdist).

  \item \textbf{Занадто широкий prior.} Якщо prior охоплює
        фізично безглузді значення ($\tau < 0$, $H_0 < 40$),
        семплер витрачає час на нефізичні регіони. Prior має
        відображати реальне фізичне знання.

  \item \textbf{Перевіряти $\hat{R}-1$, а не тільки візуально.}
        Красивий графік не гарантує збіжність. Cobaya виводить
        $\hat{R}-1$ в лог~--- він має бути $< 0.01$.

  \item \textbf{Неправильні одиниці в CAMB.} За замовчуванням
        CAMB повертає $C_\ell$, не $\mathcal{D}_\ell$. Завжди
        передавайте \texttt{CMB\_unit='muK'} і перевіряйте
        чи ділите на $2\pi/(\ell(\ell+1))$ якщо потрібно.

  \item \textbf{Порівнювати $\chi^2$, а не $\Delta\chi^2$.}
        Абсолютне значення $\chi^2$ залежить від нормування
        даних. Для порівняння двох моделей важливий
        $\Delta\chi^2 = \chi^2_{\rm model} - \chi^2_{\Lambda CDM}$
        і критерій AIC/BIC для penalізації зайвих параметрів.
\end{enumerate}

%% --------------------------------------------------------
\section{Що читати далі}
%% --------------------------------------------------------

\begin{itemize}
  \item \textbf{Документація CAMB:}
        \url{https://camb.readthedocs.io}~--- повний опис
        всіх параметрів і прикладів.
  \item \textbf{Документація Cobaya:}
        \url{https://cobaya.readthedocs.io}~--- включно з
        туторіалом по Planck likelihood.
  \item \textbf{Lesgourgues~(2011)} \cite{lesgourgues_2011}~---
        детальний опис фізики за кодом CLASS.
  \item \textbf{Lewis \& Bridle~(2002)} \cite{lewis_2002}~---
        оригінальна стаття про CosmoMC, де пояснена ідея
        MCMC для космологічних параметрів.
  \item \textbf{Planck~2018 likelihood paper}
        \cite{planck_2018_like}~--- як влаштована офіційна
        likelihood і що таке \texttt{Plik}.
\end{itemize}